% Typographic preamble for the model notes rendered as LaTeX articles.
% Loaded through pandoc's header-includes, after amsmath/booktabs/hyperref.

\usepackage{microtype}
\usepackage{mathtools}

% Headings: small-caps sections, a little air above them.
\usepackage[explicit]{titlesec}
\titleformat{\section}{\normalfont\large\bfseries}{\thesection}{1em}{#1}
\titleformat{\subsection}{\normalfont\normalsize\bfseries}{\thesubsection}{1em}{#1}
\titlespacing*{\section}{0pt}{1.6em plus .4em minus .2em}{0.7em}
\titlespacing*{\subsection}{0pt}{1.2em plus .3em minus .2em}{0.5em}

% Tables: pandoc emits booktabs longtables; give them room to breathe.
\renewcommand{\arraystretch}{1.15}
\setlength{\tabcolsep}{6pt}

% Tighter lists than the LaTeX default, which is built for standalone items.
\usepackage{enumitem}
\setlist{itemsep=0.25em, topsep=0.4em, parsep=0pt}

% Display math is frequent and often tall; the defaults leave it cramped.
\setlength{\abovedisplayskip}{10pt plus 2pt minus 2pt}
\setlength{\belowdisplayskip}{10pt plus 2pt minus 2pt}
\allowdisplaybreaks

% The combined document is a report, where each note is a part: restart the
% chapter numbering at every part so that a "§4.2" reference inside a note
% still points at the section the Markdown numbered 4.2. In the standalone
% articles there are no chapters, and nothing below applies.
\makeatletter
\@ifundefined{chapter}{}{\@addtoreset{chapter}{part}}
\makeatother
